\chapter{Parsing expression grammars}

\section{Introdução}

No capítulo anterior, estudamos os combinadores de parsing como uma
forma de implementar analisadores sintáticos descendentes recursivos.
Neste capítulo, vamos estudar um formalismo que permite a definição
de analisadores sintáticos diretamente, sem a necessidade de usar uma
gramática livre de contexto. Esse formalismo são as \textbf{Parsing
Expression Grammars} (PEGs), introduzidas por Bryan Ford em 2004.

As PEGs compartilham com as gramáticas livres de contexto (GLCs) a
ideia de descrever linguagens por meio de regras de reescrita, mas
diferem em um ponto fundamental: enquanto as GLCs usam
\textbf{escolha não-determinística} entre as alternativas de uma
produção, as PEGs usam \textbf{escolha ordenada}. Em vez de tentar
todas as alternativas em paralelo (como faz um autômato
não-determinístico), uma PEG experimenta as alternativas em ordem e
se compromete com a primeira que tiver sucesso. Essa semântica
determinística elimina a ambiguidade presente em gramáticas e
expressões regulares.

O código discutido neste capítulo está disponível nos módulos
\begin{flushleft}
  \verb|PEG|, \verb|PEG.QQ|, \verb|PEG.Examples.Arith| e
  \verb|PEG.Examples.Layout|.
\end{flushleft}

\section{Sintaxe e semântica de PEGs}

Uma PEG é composta por um conjunto de \textbf{regras de parsing},
cada uma associando um nome (não-terminal) a uma \textbf{expressão de
parsing}. As expressões de parsing são construídas a partir dos
seguintes operadores:

\begin{table}[ht]
  \centering
  \begin{tabular}{lll}
    \toprule
    Expressão & Nome & Descrição \\
    \midrule
    $\varepsilon$   & vazio              & Sempre tem sucesso sem
    consumir entrada \\
    $a$             & terminal           & Consome e casa com o símbolo $a$ \\
    $A$             & não-terminal       & Invoca a regra de nome $A$ \\
    $e_1\; e_2$     & sequência          & Tenta $e_1$ e depois $e_2$ \\
    $e_1 / e_2$     & escolha ordenada   & Tenta $e_1$; se falhar,
    tenta $e_2$ \\
    $e^*$           & repetição          & Zero ou mais repetições de
    $e$ (gulosa) \\
    $e^+$           & repetição positiva & Uma ou mais repetições de $e$ \\
    $e\,?$          & opcional           & Zero ou uma ocorrência de $e$ \\
    $!e$            & predicado negativo & Sucesso se $e$ falhar; não
    consome entrada \\
    $\& e$          & predicado positivo & Sucesso se $e$ tiver
    sucesso; não consome entrada \\
    \bottomrule
  \end{tabular}
  \caption{Operadores de PEG}
\end{table}

Uma PEG é então um conjunto de regras da forma \(A \leftarrow e\), em
que \(A\) é um não-terminal e \(e\) é uma expressão de parsing. Uma
das regras é designada como \textbf{regra inicial}.

Como exemplo, a seguinte PEG descreve a linguagem de expressões
aritméticas com precedência e associatividade já embutidas na gramática:

\[
  \begin{array}{lcl}
    E & \leftarrow & T\,(\texttt{+}\, T)^\star\\
    T & \leftarrow & F\,(*\, F)^\star\\
    F & \leftarrow & \texttt{(}\, E\, \texttt{)} \,/\, [0\text{-}9]^+
  \end{array}
\]

Comparada à GLC equivalente, a PEG não precisa de produções
recursivas à esquerda para expressar associatividade à esquerda: o
operador \((\cdot)^*\) captura diretamente a ideia de ``zero ou mais
ocorrências à direita'', que é processada iterativamente.

A semântica de uma PEG é definida em termos de uma função de parsing
que, dado uma expressão \(e\) e uma entrada \(s\), retorna ou um
\textbf{sucesso} que consiste de um par formado com a parte da
entrada consumida e o restante ainda não processado; ou uma \textbf{falha}.

A diferença crucial em relação às GLCs está na semântica da
\textbf{escolha ordenada} \(e_1 / e_2\):
Tenta-se \(e_1\) sobre a entrada \(s\). Se \(e_1\) \textbf{tem
sucesso}, o resultado é imediatamente retornado e \(e_2\) nunca é
tentada. Se \(e_1\) \textbf{falha}, então \(e_2\) é executada sobre
a entrada \(s\).

\subsection{Predicados}

Os predicados \(!e\) e \(\&\,e\) permitem analisar parte da entrada
sem necessariamente consumir parte dela. Temos os seguintes
operadores de predicado:

\begin{itemize}
  \item
    \(!e\) (\textbf{predicado negativo}): tem sucesso se e somente se
    \(e\) falha. Nenhuma entrada é consumida em nenhum caso.
  \item
    \(\&\,e\) (\textbf{predicado positivo}): tem sucesso se e somente
    se \(e\) tem sucesso. Nenhuma entrada é consumida. É equivalente
    a \(!(!e)\).
\end{itemize}

Os predicados são úteis para especificar restrições contextuais sem
avançar na entrada. Por exemplo, para reconhecer uma palavra
reservada como if sem aceitar um
identificador como iff, podemos escrever:

\[
  \texttt{if}\, !\,[a\text{-}zA\text{-}Z0\text{-}9]
\]

Isso garante que \haskell|if| só é reconhecido quando
não seguido de um caractere alfanumérico.

\subsection{Repetição}

O operador \(e^*\) é definido recursivamente como:

\[
  e^* \leftarrow e\, e^* \,/\,\varepsilon
\]

Como a escolha é ordenada, a repetição é \textbf{gulosa}: o parser
sempre tenta consumir mais uma ocorrência de \(e\) antes de aceitar o
caso vazio. Isso garante que \(e^*\) sempre consumirá o maior prefixo
possível que case com \(e\).

Ao contrário das GLCs, não há ambiguidade em \(e^*\): o resultado é
determinístico e único para qualquer entrada.

\subsection{Terminação em PEGs}

PEGs podem ser vistas como uma representação de analisadores
descendentes recursivos e, portanto, possuem limitações similares. A
terminação do processo de parsing de um PEG é garantido de terminar
se a gramática atende as seguintes restrições:

\begin{itemize}
  \item
    Não possui regras recursivas à esquerda, diretas ou indiretas;
  \item
    Toda ocorrência de repetição, \(e^*\) é tal que \(e\) não deve
    ter sucesso sem consumir símbolos da entrada.
\end{itemize}

A primeira restrição é óbvia: pois faz com que analisadores
descendentes entrem em loop. A segunda é necessária devido à
semântica gulosa do operador de repetição: ele sempre tentará casar a
expressão \(e\) antes de finalizar a repetição. Se \(e\) tiver
sucesso sem consumir entrada, então \(e^*\) pode, a cada iteração,
consumir um prefixo vazio da entrada, o que leva a não terminação.

Dizemos que uma PEG é \textbf{bem formada} se ela atende as duas
restrições mencionadas. Para PEGs bem formadas há uma garantia formal
de que o processo de parsing sempre termina.

\section{Implementação tipada de PEGs em Haskell}

A implementação de PEGs em Haskell neste capítulo segue uma abordagem
sofisticada baseada em tipos dependentes superficiais (\emph{lightweight
dependent types}): as garantias de boa-formação são verificadas
\textbf{em tempo de compilação} pelo sistema de tipos de Haskell, por
meio de GADTs, famílias de tipos e restrições de classe.

A ideia central é associar a cada expressão de parsing um
\textbf{tipo de PEG} \((\eta, F)\), onde \(\eta\) indica se a expressão pode
ter sucesso sem consumir entrada (\emph{nullable}) e \(F\) é um
sobre-conjunto dos não-terminais que podem aparecer na posição
mais à esquerda de uma derivação (\emph{head set}). Com essas
informações disponíveis no nível de tipos, o compilador Haskell rejeita
automaticamente gramáticas mal-formadas.

\subsection{Tipos de PEG}

O tipo de PEG é representado em Haskell como um tipo de dados
promovido:

\begin{minted}{haskell}
data Ty = MkTy Bool [Symbol]
\end{minted}

O primeiro componente, \haskell|Bool|, é a anulabilidade \(\eta\), e o
segundo, \haskell|[Symbol]|, é o conjunto cabeça \(F\), representado
como uma lista de nomes de não-terminais (do tipo \haskell|Symbol| de
\verb|GHC.TypeLits|). Dois seletores de família de tipos extraem os
componentes:

\begin{minted}{haskell}
type family Nullable (t :: Ty) :: Bool where
  Nullable ('MkTy n _) = n

type family First (t :: Ty) :: [Symbol] where
  First ('MkTy _ f) = f
\end{minted}

\subsection{Ambiente de tipagem}

O ambiente de tipagem \(\Gamma\) mapeia cada nome de não-terminal ao
seu tipo de PEG e ao tipo Haskell que a regra produz. No código, isso
é representado como:

\begin{minted}{haskell}
data EnvEntry = EnvEntry Ty Type
type Env = [(Symbol, EnvEntry)]
\end{minted}

O usuário declara o ambiente como uma lista de pares no nível de tipos.
Por exemplo, para a gramática de expressões aritméticas:

\begin{minted}{haskell}
type ArithEnv =
  '[ '("expr"  , 'EnvEntry ('MkTy 'False '["term","factor","number"]) Exp)
   , '("term"  , 'EnvEntry ('MkTy 'False '["factor","number"])        Exp)
   , '("factor", 'EnvEntry ('MkTy 'False '["number"])                 Exp)
   , '("number", 'EnvEntry ('MkTy 'False '[])                         Exp)
   ]
\end{minted}

Cada entrada declara o nome do não-terminal, seu tipo de PEG
\((\eta, F)\), e o tipo Haskell que a regra produz (\haskell|Exp|, neste
caso).

\subsection{O GADT \texttt{PExp}}

As expressões de parsing são representadas pelo GADT \haskell|PExp|,
indexado pelo ambiente \haskell|env|, pelo tipo de PEG \haskell|ty| e
pelo tipo Haskell do resultado \haskell|a|:

\begin{minted}{haskell}
data PExp (env :: Env) (ty :: Ty) (a :: Type) where
  Pure    :: a -> PExp env ('MkTy 'True '[]) a
  Term    :: Char -> PExp env ('MkTy 'False '[]) Char
  AnyChar :: PExp env ('MkTy 'False '[]) Char
  NT      :: (KnownSymbol s, KnownMember s env ...)
          => Name s -> PExp env (NTTy s env) (ResOf (Lookup s env))
  Seq     :: PExp env t1 (a -> b) -> PExp env t2 a
          -> PExp env (SeqTy t1 t2) b
  Choice  :: PExp env t1 a -> PExp env t2 a
          -> PExp env (ChoiceTy t1 t2) a
  Star    :: PExp env ('MkTy 'False f) a
          -> PExp env ('MkTy 'True f) [a]
  Not     :: PExp env ('MkTy n f) a -> PExp env ('MkTy 'True f) ()
  Map     :: (a -> b) -> PExp env ty a -> PExp env ty b
  Indent  :: Rel n -> PExp env ty a -> PExp env ty a
  Position :: Rel n -> PExp env ty a -> PExp env ty a
  Align   :: PExp env ty a -> PExp env ty a
\end{minted}

Cada construtor codifica a conclusão da regra de tipagem correspondente
no seu tipo de retorno:

\begin{itemize}
  \item \haskell|Pure x| sempre tem sucesso sem consumir entrada
    (\haskell|'True|) e nenhum não-terminal pode aparecer à esquerda
    (\haskell|'[]|).
  \item \haskell|Term c| consome exatamente um terminal
    (\haskell|'False|).
  \item \haskell|Star e| exige, em seu tipo, que \haskell|e| seja não
    anulável (\haskell|'MkTy 'False f|): esta é a verificação estática
    que impede laços infinitos na repetição.
  \item \haskell|Seq ef ex| combina os tipos de PEG de \haskell|ef| e
    \haskell|ex| via \haskell|SeqTy|, calculando a anulabilidade
    conjunta e a união dos conjuntos cabeça.
  \item \haskell|Choice e1 e2| combina via \haskell|ChoiceTy|,
    calculando a disjunção das anulabilidades e a união dos conjuntos
    cabeça.
\end{itemize}

Os três construtores finais (\haskell|Indent|, \haskell|Position|,
\haskell|Align|) são a extensão para parsing sensível a indentação e
serão discutidos na Seção~\ref{sec:peg-indent}.

\subsection{Gramáticas e a condição de aciclicidade}

Uma gramática completa é um valor do tipo:

\begin{minted}{haskell}
data Grammar (env :: Env) (startTy :: Ty) (startA :: Type) where
  Grammar :: Acyclic env
          => Rules env env
          -> PExp env startTy startA
          -> Grammar env startTy startA
\end{minted}

O construtor \haskell|Grammar| exige duas coisas:

\begin{enumerate}
  \item \textbf{Condição de ponto fixo}: o argumento \haskell|Rules env env|
    é uma lista heterogênea que associa cada nome do ambiente ao corpo
    de sua regra, tipado exatamente com o tipo de PEG e o tipo de
    resultado declarados no ambiente. Isso garante que as declarações
    sejam consistentes com as definições.

  \item \textbf{Aciclicidade}: a restrição \haskell|Acyclic env| é uma
    família de tipos que percorre o ambiente e verifica, para cada
    não-terminal \(i\), que \(i \notin \Gamma(i).F\). Como \(F\) é
    uma sobre-aproximação dos não-terminais que podem aparecer na
    posição mais à esquerda, essa condição rejeita tanto recursão
    direta à esquerda quanto recursão indireta. Gramáticas com recursão
    à esquerda são rejeitadas em tempo de compilação com uma mensagem
    de erro descritiva.
\end{enumerate}

\subsection{O interpretador}

A função \haskell|parse| executa uma gramática em uma entrada:

\begin{minted}{haskell}
data Result a = OK a String String | Fail

parse :: Grammar env ty a -> String -> Result a
\end{minted}

O resultado \haskell|OK v consumed rest| contém o valor produzido, a
parte da entrada consumida e o restante. \haskell|Fail| indica que a
expressão inicial não casou.

Internamente, o interpretador mantém um estado \haskell|PState|
composto pela entrada restante anotada com colunas, um conjunto de
candidatos de indentação (representado como intervalo) e uma flag de
alinhamento. A função \haskell|eval| percorre o GADT \haskell|PExp| de
forma estrutural, sem precisar de \haskell|unsafeCoerce|, pois as
testemunhas de pertinência \haskell|Member| permitem recuperar cada
corpo de regra exatamente no tipo correto.

\subsection{O quasi-quoter}

Escrever gramáticas diretamente no GADT \haskell|PExp| é preciso mas
verboso. O módulo \haskell|PEG.QQ| fornece dois quasi-quoters que
permitem usar a sintaxe concreta de Ford com rótulos e ações semânticas:

\begin{minted}{haskell}
import PEG.QQ (pegExpr, pegRules)
\end{minted}

Com \haskell|[pegRules| ... |]|, as regras são escritas na forma
\verb|nome <- expressão|, onde a expressão segue a sintaxe de Ford.
\textbf{Rótulos} (\verb|x:e|) associam nomes aos valores produzidos
pelos sub-parsers; \textbf{ações semânticas} \verb|{ expr }| constroem
o resultado usando os rótulos como variáveis em escopo:

\begin{minted}{haskell}
[pegRules|
  expr <- t:term ts:(o:[+-] u:term)* { foldl addOp t ts }
  term <- f:factor fs:(o:[*/] g:factor)*
            { foldl (\acc (op,r) -> addOp acc (op,r)) f fs }
  factor <- n:number / '(' e:expr ')' / '-' f:factor { Neg f }
  number <- ds:[0-9]+ { Lit (read ds :: Int) }
|]
\end{minted}

A sequência sem rótulos produz \haskell|()|; com um único rótulo,
produz o valor desse rótulo; com vários rótulos, produz a tupla.

Com \haskell|[pegExpr| ... |]|, uma única expressão é compilada e
pode ser usada como expressão inicial da gramática.

\section{Sensibilidade a indentação}
\label{sec:peg-indent}

A extensão para parsing sensível a indentação é baseada no trabalho de
Nestra (2017), que adiciona três operadores à sintaxe de PEGs sem
alterar o tipo de PEG do operando:

\begin{table}[ht]
  \centering
  \begin{tabular}{lll}
    \toprule
    Operador & Sintaxe concreta & Descrição \\
    \midrule
    $e^\rho$  & \verb|e^R|  & Indentação: o bloco de $e$ tem baseline
    em relação $\rho$ ao bloco externo \\
    $e_\sigma$ & \verb|e_R| & Modo de token: terminais de $e$ devem
    estar na relação $\sigma$ ao baseline \\
    $|e|$     & \verb-|e|-  & Alinhamento: o primeiro terminal de $e$
    deve estar exatamente no baseline \\
    \bottomrule
  \end{tabular}
  \caption{Operadores de indentação de PEGs}
\end{table}

As relações \(R\) disponíveis são:

\begin{itemize}
  \item \verb|>| — estritamente mais indentado (\haskell|gtR|);
  \item \verb|>=| — indentado ou alinhado (\haskell|geR|);
  \item \verb|=| — mesmo baseline (\haskell|eqR|);
  \item \verb|~| — sem restrição (\haskell|anyR|), usado para espaço
    em branco;
  \item \verb|+n| — indentado exatamente \(n\) colunas à direita.
\end{itemize}

Uma regra fundamental é que o espaço em branco deve ser consumido
sob o modo \verb|~|, para não fixar o baseline indevidamente:

\begin{minted}{haskell}
ws <- [ \t\r\n]*_~
\end{minted}

O alinhamento deve envolver o item, não o espaço em branco que o
precede: escreve-se \verb|ws |e||, nunca \verb-|ws e|-.

\section{Exemplos}

\subsection{Parser de expressões aritméticas}

O módulo \haskell|PEG.Examples.Arith| implementa um parser de expressões
que constrói uma árvore sintática abstrata. A AST é:

\begin{minted}{haskell}
data Exp
  = Lit Int
  | Neg Exp
  | Add Exp Exp | Sub Exp Exp
  | Mul Exp Exp | Div Exp Exp
\end{minted}

O ambiente de tipos declara quatro não-terminais, todos produzindo
\haskell|Exp|, e especifica os conjuntos cabeça que satisfazem a
condição de aciclicidade:

\begin{minted}{haskell}
type ArithEnv =
  '[ '("expr"  , 'EnvEntry ('MkTy 'False '["term","factor","number"]) Exp)
   , '("term"  , 'EnvEntry ('MkTy 'False '["factor","number"])        Exp)
   , '("factor", 'EnvEntry ('MkTy 'False '["number"])                 Exp)
   , '("number", 'EnvEntry ('MkTy 'False '[])                         Exp)
   ]
\end{minted}

A gramática usa o quasi-quoter para escrever as regras de forma
concisa. A associatividade à esquerda é expressa por repetição com
\haskell|foldl|:

\begin{minted}{haskell}
arith :: Grammar ArithEnv _ Exp
arith =
  Grammar
    [pegRules|
       expr   <- t:term ts:(o:[+-] u:term)* { foldl addOp t ts }
       term   <- f:factor fs:(o:[*/] g:factor)*
                   { foldl (\acc (op,r) -> addOp acc (op,r)) f fs }
       factor <- n:number
               / '(' e:expr ')'
               / '-' f:factor { Neg f }
       number <- ds:[0-9]+ { Lit (read ds :: Int) }
    |]
    (nt @"expr")
\end{minted}

Para executar o parser:

\begin{minted}{haskell}
ghci> parse arith "1+2*3"
OK (Add (Lit 1) (Mul (Lit 2) (Lit 3))) "1+2*3" ""

ghci> parse arith "(1+2)*3"
OK (Mul (Add (Lit 1) (Lit 2)) (Lit 3)) "(1+2)*3" ""

ghci> parse arith "10-2-3"
OK (Sub (Sub (Lit 10) (Lit 2)) (Lit 3)) "10-2-3" ""
\end{minted}

Note que \haskell|10-2-3| é associado à esquerda, como esperado: o
\haskell|foldl| no quasi-quoter cuida disso.

\subsection{Blocos ao estilo Haskell}

O módulo \haskell|PEG.Examples.Layout| implementa um parser para
expressões \haskell|do| com dois modos de layout: indentação (como em
Haskell) e colchetes explícitos (como em Haskell com a extensão
\emph{NoImplicitPrelude}).

A AST é simples:

\begin{minted}{haskell}
data DoStmt = Atom String | Nested [DoStmt]
\end{minted}

O ambiente de tipos:

\begin{minted}{haskell}
type DoEnv =
  '[ '("doexp" , 'EnvEntry ('MkTy 'False '[])                               [DoStmt])
   , '("istmts", 'EnvEntry ('MkTy 'False '["ws","stmt","doexp","name"])     [DoStmt])
   , '("stmts" , 'EnvEntry ('MkTy 'False '["ws"])                          [DoStmt])
   , '("stmt"  , 'EnvEntry ('MkTy 'False '["doexp","name"])               DoStmt)
   , '("name"  , 'EnvEntry ('MkTy 'False '[])                             String)
   , '("ws"    , 'EnvEntry ('MkTy 'True  '[])                             ())
   ]
\end{minted}

A gramática usa os três operadores de indentação:

\begin{minted}{haskell}
doExp :: Grammar DoEnv _ [DoStmt]
doExp =
  Grammar
    [pegRules|
       doexp  <- "do" b:(i:istmts / j:stmts)
       istmts <- ss:(ws st:|s:stmt|)+^>
       stmts  <- r:(ws '{' ws s:stmt ss:(ws ';' ws t:stmt)* ws '}' { s:ss })^~
       stmt   <- d:doexp { Nested d } / n:name { Atom n }
       name   <- cs:[a-z]+
       ws     <- [ \t\r\n]*_~
    |]
    [pegExpr| ws d:doexp ws !. |]
\end{minted}

Na regra \texttt{istmts}, a expressão
\texttt{(ws st:{\textbar}s:stmt{\textbar})+\^{}>}
constrói uma lista de declarações alinhadas (\texttt{{\textbar}s:stmt{\textbar}})
em um bloco estritamente mais indentado (\texttt{\^{}>}) do que o bloco externo.
Na regra \texttt{stmts}, o grupo \texttt{\^{}\~{}} usa a relação irrestrita,
que é o que os colchetes explícitos exigem.

Para executar com opções de layout:

\begin{minted}{haskell}
ghci> parseWith layoutOpts doExp "do a\n   b\n   c\n"
OK [Atom "a", Atom "b", Atom "c"] ...

ghci> parseWith layoutOpts doExp "do { a; b; c }\n"
OK [Atom "a", Atom "b", Atom "c"] ...

ghci> parseWith layoutOpts doExp "do a\n   do x\n      y\n   b\n"
OK [Atom "a", Nested [Atom "x", Atom "y"], Atom "b"] ...
\end{minted}

O parser aceita tanto o layout por indentação quanto os colchetes
explícitos, e os dois modos podem ser aninhados arbitrariamente.

\section{PEGs e GLCs: principais diferenças}

As PEGs e as GLCs são formalismos distintos para descrever
linguagens. As principais diferenças práticas são:

\textbf{Ambiguidade.} Toda PEG é, por definição, não-ambígua: a
escolha ordenada garante que exista no máximo um resultado para
qualquer entrada. As GLCs podem ser ambíguas, e resolver ambiguidades
muitas vezes exige reformular a gramática ou usar anotações de
precedência externas (como os mecanismos
  \verb|%left|/\verb|%right| do yacc/bison).

  \textbf{Recursão à esquerda.} GLCs admitem recursão à esquerda
  naturalmente; por exemplo, \(E \to E + T\) é uma produção válida e
  útil. Parsers PEG descendentes não suportam recursão à esquerda
  direta. A solução é reescrever usando repetição \((\cdot)^*\) ou
  \haskell|foldl|.

  \textbf{Poder expressivo.} PEGs reconhecem linguagens que as GLCs não
  conseguem (por exemplo, \(\{a^n b^n c^n \mid n \geq 0\}\), usando
  predicados) e há conjectura de que algumas linguagens livres de
  contexto não podem ser descritas por PEGs, embora isso não esteja
  formalmente provado.

  \textbf{Sensibilidade a indentação.} A extensão de Nestra permite
  que PEGs descrevam linguagens sensíveis a indentação, como Python e
  Haskell, de forma declarativa, sem pré-processamento léxico.

  \section{Conclusão}

  Este capítulo apresentou as Parsing Expression Grammars como um
  formalismo determinístico para especificar analisadores sintáticos.
  Apresentamos uma implementação tipada em Haskell que usa GADTs e
  famílias de tipos para verificar, em tempo de compilação, as
  condições de boa-formação: anulabilidade não-trivial nos corpos de
  repetição e aciclicidade do ambiente de tipagem.

  Os quasi-quoters \haskell|pegRules| e \haskell|pegExpr| permitem
  escrever gramáticas na sintaxe concreta de Ford com rótulos e ações
  semânticas, sem abrir mão das garantias estáticas. A extensão para
  indentação de Nestra adiciona três operadores (\texttt{\^{}},
  \texttt{\_}, \texttt{|e|}) que permitem descrever linguagens
  sensíveis a layout de forma puramente declarativa.
